% p180_1; source PDF page 180; manually verified editable table.
\begin{tikzpicture}[x=1cm,y=-1cm,inner sep=0pt,outer sep=0pt]
\path[use as bounding box] (0,0) rectangle (16.000,6.400);
\draw[line width=.4pt] (0.000,0.000) rectangle (2.400,0.800);
\node[align=center,text width=2.140cm,font={\linespread{1}\fontsize{10}{13.50}\selectfont\bfseries}] at (1.200,0.400) {基本原理};
\draw[line width=.4pt] (2.400,0.000) rectangle (6.500,0.800);
\node[align=center,text width=3.840cm,font={\linespread{1}\fontsize{10}{13.50}\selectfont\bfseries}] at (4.450,0.400) {处理办法};
\draw[line width=.4pt] (6.500,0.000) rectangle (16.000,0.800);
\node[align=center,text width=9.240cm,font={\linespread{1}\fontsize{10}{13.50}\selectfont\bfseries}] at (11.250,0.400) {范例};
\draw[line width=.4pt] (0.000,0.800) rectangle (2.400,6.400);
\node[align=left,text width=2.140cm,font={\linespread{1}\fontsize{10}{13.50}\selectfont}] at (1.200,3.600) {为了保护法益，禁止遗漏处罚（评价）};
\draw[line width=.4pt] (2.400,0.800) rectangle (6.500,4.600);
\node[align=left,text width=3.840cm,font={\linespread{1}\fontsize{10}{13.50}\selectfont}] at (4.450,2.700) {原则上，数个行为，侵犯数个法益，触犯数个罪名，就应数罪并罚，否则会遗漏处罚其中一个行为制造的法益侵害事实};
\draw[line width=.4pt] (6.500,0.800) rectangle (16.000,4.600);
\node[align=left,text width=9.240cm,font={\linespread{1}\fontsize{10}{13.50}\selectfont}] at (11.250,2.700) {甲暴力抢劫乙，乙逃跑，途中不慎钱包落地，乙不知，跑远了。甲捡走钱包。甲的暴力行为构成抢劫罪未遂，捡走钱包构成侵占罪既遂，应数罪并罚。若用抢劫罪吸收侵占罪，则遗漏处罚拿走钱包这一法益侵害结果。若只定侵占罪，则遗漏评价了暴力行为};
\draw[line width=.4pt] (2.400,4.600) rectangle (6.500,6.400);
\node[align=left,text width=3.840cm,font={\linespread{1}\fontsize{10}{13.50}\selectfont}] at (4.450,5.500) {例外的，作一罪处理};
\draw[line width=.4pt] (6.500,4.600) rectangle (16.000,6.400);
\node[align=left,text width=9.240cm,font={\linespread{1}\fontsize{10}{13.50}\selectfont}] at (11.250,5.500) {作一罪处理，要么具有法条依据，例如结合犯；要么具有充分的理论依据，例如吸收犯、牵连犯};
\end{tikzpicture}
